\section{Hurwitz endpoints and all Stieltjes orders}
\label{sec:hurwitz}

The new coefficient densities naturally supply the polynomial family
\begin{equation}\label{eq:Wkernel}
 W_{b,n}(v)=(1-v^2)^3(1+3v^2)^b(v+v^3)^{n-2},
 \qquad b\ge0,\ n\ge2.
\end{equation}
Multiplication by $v^k$ gives all coefficient moments in
\eqref{eq:coefficient-integral}. Every such polynomial and its first two
derivatives vanish at $v=\pm1$. This creates a finite endpoint calculus
starting with the third derivative, uniformly in $b,n,k$.

\subsection{The polynomial endpoint identity}
Let $P$ be any polynomial of degree $d$ and define
\[
 M_P(s;a)=\frac12\int_{-1}^1P(v)\zeta(s,a+v)\dd v,
 \qquad \mu_P=\avg{P(v)}.
\]
Use $\fall{s-1}{j+1}=(s-1)(s-2)\cdots(s-j-1)$.

\begin{proposition}[Finite Hurwitz endpoint formula]\label{prop:Hurwitz}
For $\Re a>1$, as an identity of meromorphic functions of $s$,
\begin{equation}\label{eq:Hurwitz-endpoint}
 M_P(s;a)=-\frac12\sum_{j=0}^d
 \frac{\eval{P^{(j)}(v)\zeta(s-j-1,a+v)}}{\fall{s-1}{j+1}}.
\end{equation}
Its only possible pole is at $s=1$, with residue $\mu_P$.
All other apparent singularities of the right side are removable after
summation. If $P,P',P''$ vanish at both endpoints, the sum starts at
$j=3$.
\end{proposition}
\begin{proof}
The classical argument derivative
$\partial_a\zeta(s-1,a)=-(s-1)\zeta(s,a)$
\cite{DLMF25} gives the first integration by parts. Repeat until
$P^{(d+1)}=0$. Each endpoint term has the minus sign shown in
\eqref{eq:Hurwitz-endpoint}. This proves the formula away from the
exceptional integers. On the compact interval and locally uniformly for
$\Re a>1$, the integrand is analytic in $s$ except for its simple
pole at one. Its residue integrates to $\mu_P$. Meromorphic continuation
therefore proves the endpoint identity and all stated cancellations.
\end{proof}

This proposition is the standard polynomial integration-by-parts
mechanism. The specialization to $W_{b,n}$ is what connects the newly
computed Gamma normal coefficients to spectral special functions.
No endpoint regularization in the variable $v$ is needed: the integrals
are ordinary compact-interval integrals. Laurent coefficients in $s$
are a different operation.

\subsection{Finite harmonic coefficients for every Stieltjes order}
For $j\ge0$ define $\eta_{j,\ell}$ by
\begin{equation}\label{eq:eta}
 \prod_{r=1}^j(1-\epsilon/r)^{-1}
 =\sum_{\ell\ge0}\eta_{j,\ell}\epsilon^\ell
 =\exp\left(\sum_{r\ge1}\frac{H_j^{(r)}}r\epsilon^r\right).
\end{equation}
The empty product at $j=0$ is one. In particular,
\[
 \eta_{j,0}=1,\qquad \eta_{j,1}=H_j,\qquad
 \eta_{j,2}=\frac{H_j^2+H_j^{(2)}}2.
\]
These are finite polynomials in generalized harmonic numbers at each
fixed $\ell$.

\begin{theorem}[All-order Stieltjes averages]\label{thm:Stieltjes}
For $q\ge0$ and $\Re a>1$,
\begin{equation}\label{eq:Stieltjes-endpoint}
\begin{split}
 \avg{P(v)\gamma_q(a+v)}
 ={}&\frac{(-1)^{q+1}q!}{2}
 \sum_{j=0}^d\frac{(-1)^j}{j!}\\[-2pt]
 &\quad\times\eval{P^{(j)}(v)
 \sum_{r=0}^{q+1}\eta_{j,q+1-r}
       \frac{\zeta^{(r)}(-j,a+v)}{r!}}.
\end{split}
\end{equation}
The formula is finite at every $q$ and applies in particular to every
polynomial $v^kW_{b,n}$.
\end{theorem}
\begin{proof}
At $s=1+\epsilon$,
\begin{equation}\label{eq:falling-laurent}
 \frac1{\fall{\epsilon}{j+1}}
 =\frac{(-1)^j}{j!\epsilon}
       \sum_{\ell\ge0}\eta_{j,\ell}\epsilon^\ell.
\end{equation}
Expand $\zeta(-j+\epsilon,a+v)$ in its order variable and insert
\eqref{eq:falling-laurent} into \eqref{eq:Hurwitz-endpoint}. Coefficient
$\epsilon^q$ is the right side of \eqref{eq:Stieltjes-endpoint}
multiplied by $(-1)^q/q!$. On the integral side the same coefficient is
$(-1)^q\avg{P\gamma_q}/q!$ by
\eqref{eq:stieltjes-convention}. Uniformity on compact parameter sets
justifies coefficient extraction under the integral. The coefficient of
$\epsilon^{-1}$ is the separately retained residue $\mu_P$; it must not
be discarded before multiplying the expansions.
\end{proof}

\begin{corollary}[Anchored primitives and argument derivatives]
\label{cor:Stieltjes-calculus}
For $a_0,a$ in the half-plane $\Re a>1$, with the integral taken along
any path in that half-plane,
\begin{equation}\label{eq:Stieltjes-primitive}
 \int_{a_0}^a\avg{P(v)\gamma_q(c+v)}\dd c
 =\left.\frac{(-1)^{q+1}}{q+1}\partial_s^{q+1}
       \bigl(M_P(s;a)-M_P(s;a_0)\bigr)\right|_{s=0}.
\end{equation}
For $r\ge1$, let $e_j^{[r]}$ be the elementary symmetric polynomial of
degree $j$ in $1,1/2,\ldots,1/r$. Then
\begin{equation}\label{eq:Stieltjes-argument-derivative}
\begin{split}
 \partial_a^r\avg{P(v)\gamma_q(a+v)}
 ={}&(-1)^{q+r}q!r!
 \sum_{j=0}^{\min(q,r)}e_j^{[r]}
 \frac{\left.\partial_s^{q-j}M_P(s;a)\right|_{s=r+1}}{(q-j)!}.
\end{split}
\end{equation}
Here the derivative $\partial_s$ on $M_P(r+1;a)$ is evaluated in its
first argument. Both formulas reduce to finite endpoint expressions by
Proposition~\ref{prop:Hurwitz} and its removable limits.
\end{corollary}
\begin{proof}
From $\partial_a\zeta(s,a)=-s\zeta(s+1,a)$ and the Laurent expansion,
\[
 \partial_a\zeta^{(q+1)}(0,a)
 =(-1)^{q+1}(q+1)\gamma_q(a).
\]
Integrate in $a$ and average against $P$ to obtain
\eqref{eq:Stieltjes-primitive}; the displayed subtraction fixes the
anchor exactly. For the derivative formula, use
\[
 \partial_a^r\zeta(s,a)=(-1)^r(s)_r\zeta(s+r,a),\qquad
 (1+\epsilon)_r=r!\sum_{j=0}^re_j^{[r]}\epsilon^j,
\]
and compare coefficients at $s=1+\epsilon$.
\end{proof}

For $q=0$, the primitive is the weighted log-Gamma difference because
$\zeta'(0,a)=\log\Gamma(a)-\tfrac12\log(2\pi)$
\cite{DLMF25}. Formula~\eqref{eq:Stieltjes-argument-derivative} at
$q=0$ recovers the weighted polygamma integrals, using
$\psi^{(r)}(a)=(-1)^{r+1}r!\zeta(r+1,a)$.

\subsection{Explicit digamma and Stieltjes evaluations}
\begin{corollary}[Four concrete evaluations]\label{cor:four-integrals}
All integrals below are ordinary integrals over $[-1,1]$:
\begin{align}
 \int_{-1}^1P_3(v)\psi(2+v)\dd v
 &=16\zeta'(-3)+12\zeta'(-5)+\frac27,
 \label{eq:I-psi-base}\\
 \int_{-1}^1P_3(v)(1+3v^2)\psi(2+v)\dd v
 &=64\zeta'(-3)+240\zeta'(-5)+48\zeta'(-7)
       +\frac{547}{1800},\label{eq:I-psi-mixed}\\
 \int_{-1}^1P_3(v)(v+v^3)\psi(2+v)\dd v
 &=\frac{32}{105}-\frac{84\zeta(5)}{\pi^4}
   +\frac{1575\zeta(7)}{2\pi^6}-\frac{2835\zeta(9)}{2\pi^8},
 \label{eq:I-psi-odd}\\
 \int_{-1}^1P_3(v)\gamma_1(2+v)\dd v
 &=8\zeta''(-3)+6\zeta''(-5)+\frac{88}{3}\zeta'(-3)
   +\frac{137}{5}\zeta'(-5)\notag\\
 &\hspace{1em}-\frac{16}{15}\log2+\frac{521}{900}.
 \label{eq:I-stieltjes}
\end{align}
\end{corollary}
\begin{proof}
At $a=2$ the endpoint arguments are $1$ and $3$. Use the exact shift
\begin{equation}\label{eq:zeta-three}
 \zeta(s,3)=\zeta(s)-1-2^{-s}.
\end{equation}
For $q=0$, Theorem~\ref{thm:Stieltjes} and $\gamma_0=-\psi$ give
\begin{equation}\label{eq:psi-endpoint}
 \int_{-1}^1P(v)\psi(2+v)\dd v
 =\sum_{j=0}^d\frac{(-1)^j}{j!}
 \eval{P^{(j)}(v)\bigl(\zeta'(-j,2+v)+H_j\zeta(-j,2+v)\bigr)}.
\end{equation}
The nonzero derivatives at $v=1$ needed for the three polynomials are
listed below; their derivatives at $-1$ follow from parity.
\begin{center}\small
\begin{tabular}{r|rrr}
\toprule
$j$ & $P_3^{(j)}(1)$ & $(P_3(1+3v^2))^{(j)}(1)$
    & $(P_3(v+v^3))^{(j)}(1)$\\
\midrule
3 & $-48$ & $-192$ & $-96$\\
4 & $-288$ & $-2304$ & $-1344$\\
5 & $-720$ & $-14400$ & $-10080$\\
6 & $-720$ & $-54720$ & $-50400$\\
7 & $0$ & $-120960$ & $-171360$\\
8 & $0$ & $-120960$ & $-362880$\\
9 & $0$ & $0$ & $-362880$\\
\bottomrule
\end{tabular}
\end{center}
Substituting these finite lists into \eqref{eq:psi-endpoint} and using
\eqref{eq:zeta-three} proves \eqref{eq:I-psi-base} and
\eqref{eq:I-psi-mixed}. The third polynomial first gives
\[
 -112\zeta'(-4)-140\zeta'(-6)-18\zeta'(-8)+\frac{32}{105}.
\]
The zeta functional equation implies
$\zeta'(-2r)=(-1)^r(2r)!\zeta(2r+1)/(2(2\pi)^{2r})$;
substitution gives \eqref{eq:I-psi-odd}.

Finally use $q=1$ in \eqref{eq:Stieltjes-endpoint}. Its inner sum is
\[
 \tfrac12\zeta''(-j,2+v)+H_j\zeta'(-j,2+v)
 +\tfrac12(H_j^2+H_j^{(2)})\zeta(-j,2+v).
\]
With the first derivative column and \eqref{eq:zeta-three}, collecting
terms yields \eqref{eq:I-stieltjes}. These calculations require no
numerical recognition of constants.
\end{proof}

The second and third integrals use precisely the first mixed square and
nonlinear cubic coefficient kernels, apart from their fixed scaling
constants. The theorem thus links the algebraic coefficients to explicit
zeta and Stieltjes periods. It does not assert that these particular
low-degree evaluations have no earlier appearance in the literature.
